\section{Data}
\label{sec:data}

\subsection{Instrument, feeds, and the bookTicker discontinuation constraint}
\label{sec:data-feeds}

The primary instrument is \texttt{BTCUSDT} on Binance USD-margined
perpetual futures (\texttt{futures/um/daily/}). Two raw public feeds are
used: unaggregated \texttt{trades} (individual execution ticks, not the
\texttt{aggTrades} feed, which collapses simultaneous fills and would
destroy the clustering signal the Hawkes process is meant to measure) and
\texttt{bookTicker} (L1 best-bid/best-ask updates). Both feeds are ordered
on \texttt{transaction\_time} (\texttt{T}), not \texttt{event\_time}
(\texttt{E}), since the latter includes Binance's own engine push
latency and is not a reliable causal ordering key.

The data window is constrained by a permanent upstream limitation:
Binance stopped publishing \texttt{futures/um/daily/bookTicker/BTCUSDT}
after \textbf{2024-03-30}. This is documented in the
\texttt{binance-public-data} public GitHub repository (Issue \#372, filed
November 2024) and remained open and unresolved as of the most recent
check (August 2026, 20+ months after filing) -- it is treated in this
paper as a permanent, not transient, constraint, and rules out any
attempt to use a more recent window. The \texttt{trades} feed itself
remains continuous and current; only \texttt{bookTicker} is affected.
Separately, Binance's own archive is documented (Issue \#305) to have
shipped BTC/ETH \texttt{bookTicker} daily files \emph{out of sequence}
during Jan--Mar 2024, so no assumption of pre-sorted archive files is
made anywhere in the pipeline; all downstream processing sorts explicitly
by \texttt{(transaction\_time, update\_id)} before any diff-based
computation.

\subsection{The 38-day primary window}
\label{sec:data-window}

An initial full-coverage audit surfaced severe raw ordering disorder in
\texttt{bookTicker} between 2024-02-04 and 2024-02-21 -- on some days in
this range, up to $\sim$40\% of a day's rows exhibited a
\texttt{transaction\_time} decrease relative to the previous row (e.g.\
18{,}728{,}015 of 46{,}396{,}761 rows on 2024-02-12), a scale of disorder
far beyond isolated same-millisecond ties. An independent re-audit
confirmed that raw ordering disorder ends precisely at 2024-02-21 --
2024-02-22 onward shows zero raw \texttt{transaction\_time} decreases.
The primary study window is therefore fixed to
\[
  \textbf{2024-02-22 -- 2024-03-30 (38 calendar days)},
\]
the last available continuous window before the \texttt{bookTicker}
discontinuation and after the resolved disorder period. An independent
re-audit of this exact window (38 days, $875{,}126{,}040$ raw
\texttt{bookTicker} rows) found \emph{zero} adjacent
\texttt{transaction\_time} decreases, zero same-millisecond
\texttt{update\_id} regressions, and zero invalid-day rows across every
one of the 38 days. Two days retain known, small activity gaps --
2024-02-28 (largest interval 17:32:13--18:15:20 UTC, 2{,}587 seconds
inactive) and 2024-03-08 (largest interval 16:03:08--17:09:20 UTC,
3{,}972 seconds inactive) -- totaling roughly two hours ten minutes of
real inactivity across the entire 38-day window; every other day is
100.000\% or near-100.000\% active on a 1-second-bin coverage measure.
The excluded 2024-02-04--02-21 disorder window is discussed further as a
scope limitation in Section~\ref{sec:limitations}.

Raw input over the primary window totals $196{,}875{,}099$
\texttt{trades} rows and $875{,}126{,}040$ \texttt{bookTicker} rows
($1{,}072{,}001{,}139$ raw rows combined). After causal alignment and
classification into the 10-event taxonomy of
Section~\ref{sec:data-taxonomy} -- which collapses redundant,
same-direction consecutive \texttt{bookTicker} updates into single typed
transitions -- the aligned event stream contains
\[
  \mathbf{1{,}053{,}434{,}969}\ \text{typed events}: \quad 196{,}875{,}099\ \text{trades} + 856{,}559{,}870\ \text{emitted quote events},
\]
with zero unclassified quote transitions on every one of the 38 days.
This just over one billion typed events is the number referred to
throughout the paper as the primary window's event volume.

\subsection{The 10-event microstructure taxonomy}
\label{sec:data-taxonomy}

Every raw quote or trade transition is mapped to exactly one of ten
mutually exclusive event types, with residual cancellation volume at
price level $p$ defined as
\[
  V_{\text{cancel}}(p) = \max\!\big(0,\ -\Delta Q(p) - V_{\text{trade}}(p)\big).
\]

\begin{center}
\begin{tabular}{@{}cll@{}}
\toprule
\# & Event type & Definition \\
\midrule
1 & \texttt{bid\_qty\_up}    & bid price unchanged, $\Delta Q_b > 0$ \\
2 & \texttt{bid\_qty\_down}  & bid price unchanged, $V_{\text{cancel}}(P_b) > 0$ \\
3 & \texttt{bid\_price\_up}   & bid price increases (spread narrowing / price improvement) \\
4 & \texttt{bid\_price\_down} & bid price decreases (queue depletion / spread widening) \\
5 & \texttt{ask\_qty\_up}    & ask price unchanged, $\Delta Q_a > 0$ \\
6 & \texttt{ask\_qty\_down}  & ask price unchanged, $V_{\text{cancel}}(P_a) > 0$ \\
7 & \texttt{ask\_price\_up}   & ask price increases (queue depletion / spread widening) \\
8 & \texttt{ask\_price\_down} & ask price decreases (spread narrowing / price improvement) \\
9 & \texttt{trade\_buy}      & aggressive buyer-initiated fill \\
10 & \texttt{trade\_sell}     & aggressive seller-initiated fill \\
\bottomrule
\end{tabular}
\end{center}

Table~\ref{tab:event-counts} gives the exact event-type counts over the
full 38-day aligned stream.

\begin{table}[htbp]
\centering
\caption{Event-type counts, 38-day aligned stream (2024-02-22 -- 2024-03-30).}
\label{tab:event-counts}
\begin{tabular}{@{}lr@{}}
\toprule
Event type & Count \\
\midrule
\texttt{bid\_qty\_up}    & 258{,}877{,}869 \\
\texttt{bid\_qty\_down}  & 151{,}038{,}105 \\
\texttt{bid\_price\_up}   & 10{,}624{,}783 \\
\texttt{bid\_price\_down} & 8{,}142{,}038 \\
\texttt{ask\_qty\_up}    & 258{,}664{,}686 \\
\texttt{ask\_qty\_down}  & 150{,}541{,}519 \\
\texttt{ask\_price\_up}   & 8{,}303{,}725 \\
\texttt{ask\_price\_down} & 10{,}367{,}145 \\
\texttt{trade\_buy}      & 98{,}416{,}666 \\
\texttt{trade\_sell}     & 98{,}458{,}433 \\
\midrule
\textbf{Total} & \textbf{1{,}053{,}434{,}969} \\
\bottomrule
\end{tabular}
\end{table}

\subsection{Timestamp alignment and the tie-break causality sensitivity check}
\label{sec:data-tiebreak}

Sub-millisecond ties within a single feed are broken by a deterministic
micro-jitter; within \texttt{bookTicker}, exact-\texttt{transaction\_time}
ties are further broken by \texttt{update\_id} before falling back to
arrival order. The default cross-feed ordering convention on an
exact-millisecond tie between the two feeds is \texttt{trades\_first}
(trade ordered before a same-millisecond \texttt{bookTicker} update).
Because this convention is an arbitrary modeling choice rather than a
directly observable fact about true event order, every Hawkes cross-excitation
result is subjected to a mandatory causality sensitivity check: the
Primary Endpoint and the fitted branching-ratio matrix
$\boldsymbol{\Gamma}$ are recomputed under the inverted rule,
\texttt{book\_first}, and compared.

On the sample day used for initial data inspection (2024-03-30),
exact-millisecond trade/\texttt{bookTicker} collisions accounted for
7.99\% of all quote timestamps (372{,}815 collisions that day), confirming
the sensitivity check is not a minor edge case. At $N=3$d on the 38-day
window, the initial (subsequently found to be confounded, see below)
sensitivity run found $M_2$ PR-AUC of $0.6805835954$ under
\texttt{book\_first} versus $0.6923842563$ under \texttt{trades\_first} --
a difference of $-0.01180066091$ -- while $M_1$ and $M_0$ were essentially
unaffected ($\Delta M_1 = +0.0008535961$, $\Delta M_0 = 0$). The
branching-ratio matrix $\boldsymbol{\Gamma}$ itself changed far more than
the downstream PR-AUC: $43.67\%$ relative Frobenius-norm change, with
$89.22\%$ of the squared change concentrated in the quote$\to$trade
block -- physically, in the same-millisecond touch-price/trade pairs
where the tie rule directly determines whether a touch-price move or the
paired trade is treated as causally prior.

A follow-up investigation isolated this effect more cleanly by holding
the event multiset exactly fixed under an ordering-only inversion (rather
than re-running the full aligner, which also changes trade-volume
attribution at qty-tie edges and confounds ordering with event-typing
effects). Under this cleaner order-only design, the pure tie-order effect
on $\boldsymbol{\Gamma}$ was found to be \emph{larger}, not smaller, than
the original confounded estimate ($61.9\%$ relative Frobenius change,
unmasked), concentrated almost entirely (95\% of squared change) in an
``ignition'' artifact: trades acquire large, spurious self-exciting
coefficients from touch-price events that previously sat strictly before
them at the same timestamp. Masking same-millisecond touch-price/trade
ties out of the Hawkes calibration (\texttt{mask\_touch\_trade}) reduces
this to an $8.4\%$ residual, itself attributable to a disclosed,
narrower and unmasked class of same-millisecond quantity-only quote/trade
ties (6.35\% of all events) whose effect propagates into the
quote$\to$quote block via ordinary multicollinearity in the joint
per-target regression, rather than a defect in the masking logic itself.
Across every configuration tested, the downstream PR-AUC sensitivity of
the pure order-only, masked comparison shrank to $-0.0002365$, roughly
$50\times$ smaller than the original confounded $-0.0118$ figure and
comparable in magnitude to this project's own documented run-to-run
numerical noise floor (Section~\ref{sec:limitations}). The full
investigation, decision log, and its extension to the cross-asset
transfer setting are summarized further in
Sections~\ref{sec:results-transfer} and~\ref{sec:limitations}.

\subsection{Fixed test-week representativeness}
\label{sec:data-testweek}

The learning-curve test set (Section~\ref{sec:methodology}) is fixed to
the final calendar week of the primary window, 2024-03-24 -- 2024-03-30
(a full Sunday-through-Saturday week, guaranteeing all weekdays, one
weekend, and three full 8-hour funding cycles per day regardless of
anchor point). This choice was checked, not assumed: every daily metric
in the test week (trade row count, quote notional, realized volatility,
\texttt{bookTicker} row count, aligned event count) was compared against
the pre-test-week distribution (2024-02-22 -- 2024-03-23) using a robust
$z$-score, $z = (x - \mathrm{median})/(1.4826\,\mathrm{MAD})$, with an
ex-ante flag threshold of $|z| > 3.5$. No test-week observation was
flagged (\texttt{test\_week\_anomalous=false}); the largest absolute
robust $z$-score observed was $2.424538$, for aligned-event count on
2024-03-30.
